\documentclass[
  % -------------------------
  % Curso
  % -------------------------
  coursetitle={Cómputo Nube y Tecnologías Emergentes},
  coursehours={96},
  coursedate={Verano 2026},
  doctype={Guía del curso},
  otherlangs={english,french},
  % -------------------------
  % Autor
  % -------------------------
  authorname={Lázaro},
  authorsurname={Bustio-Martínez},
  authortitle={PhD},
  role={Profesor Investigador},
  email={lbustio@inaoe.mx},
  phone={5559504000},
  ext={7459},
  % -------------------------
  % Institución
  % -------------------------
  institution={Instituto Nacional de Astrofísica, Óptica y Electrónica (INAOE)},
  department={Maestría en Ciencias en Tecnologías de Seguridad},
  %unit={Coordinación de Ciberseguridad},
  address={Luis Enrique Erro 1, Sta. Ma. Tonantzintla, San Andrés Cholula, CP 72840, Puebla, México},
  % -------------------------
  % Márgenes
  % -------------------------
  margin={0.5in}
]{../../lbustio_lecture_notes}

\usepackage{longtable}
\usepackage{array}
\usepackage{url}
\usepackage{enumitem}

\begin{document}

\IberoCover

\section{Descripción general}
\label{sec:descripcion}

Este curso introduce los fundamentos del cómputo en la nube y de diversas tecnologías emergentes con aplicaciones en ciberseguridad. Se estudian los modelos de servicio y despliegue de la nube, la gestión distribuida de datos, los principales mecanismos de seguridad en entornos cloud y el uso de plataformas de inteligencia artificial y aprendizaje automático. Asimismo, se analizan los fundamentos de blockchain, los contratos inteligentes, los criptoactivos y algunas de sus aplicaciones más relevantes.

El curso combina bases conceptuales con actividades prácticas orientadas al diseño, implementación y evaluación de soluciones tecnológicas seguras apoyadas en infraestructura cloud.

\section{Objetivos}
\label{sec:objetivos}

A continuación se presentan los objetivos general y específicos del curso.

\subsection{Objetivo general}
\label{sec:objetivo_general}

Proporcionar al estudiante los conocimientos teóricos y prácticos necesarios para comprender, analizar y aplicar los fundamentos del cómputo en la nube y de tecnologías emergentes, tales como inteligencia artificial y blockchain, con el fin de diseñar e implementar soluciones seguras para problemas relevantes en el ámbito de la ciberseguridad.

\subsection{Objetivos específicos}
\label{sec:objetivos_especificos}
\begin{itemize}
  \item Comprender los modelos de servicio, despliegue y operación del cómputo en la nube, considerando elasticidad, virtualización, costos, conectividad y riesgos de seguridad.
  \item Desplegar y administrar infraestructura básica en Google Cloud Platform, incluyendo máquinas virtuales, almacenamiento, acceso remoto, reglas de firewall y servicios web funcionales.
  \item Construir pipelines de ingestión, búsqueda y visualización de eventos mediante Elasticsearch y Kibana sobre infraestructura cloud real.
  \item Recolectar, centralizar y analizar telemetría y logs de seguridad para identificar patrones, actividad sospechosa y limitaciones del monitoreo en entornos cloud.
  \item Aplicar servicios cloud de inteligencia artificial y aprendizaje automático, como BigQuery ML y Vertex AI, para explorar datasets de eventos, detectar anomalías e interpretar resultados.
  \item Desplegar e interactuar con una blockchain educativa privada para comprender bloques, hashes, transacciones, integridad, consenso simplificado y aplicaciones básicas en ciberseguridad.
  \item Documentar de forma reproducible las arquitecturas, procedimientos, problemas encontrados, decisiones técnicas, evidencias de funcionamiento y reflexiones críticas de cada práctica.
\end{itemize}

\section{Estructura del curso}
\label{sec:estrucrtura}
El curso está organizado en unidades progresivas que combinan fundamentos conceptuales, análisis de arquitecturas, demostraciones técnicas y prácticas guiadas sobre infraestructura cloud. La parte teórica introduce los principios que explican el funcionamiento de los sistemas distribuidos modernos: virtualización, contenedores, redes definidas por software, almacenamiento administrado, elasticidad, alta disponibilidad, observabilidad, control de acceso, gestión de identidades, cifrado, automatización, servicios de inteligencia artificial y mecanismos básicos de blockchain. La parte práctica busca que el estudiante despliegue, configure, documente y evalúe soluciones reales, identificando tanto sus capacidades como sus riesgos de seguridad.

Las unidades del curso se articulan alrededor de cuatro ejes. El primero aborda los fundamentos del cómputo en la nube, incluyendo modelos de servicio, modelos de despliegue, regiones, zonas, escalabilidad, costos y responsabilidades compartidas de seguridad. El segundo se centra en infraestructura y servicios administrados: máquinas virtuales, almacenamiento, redes, reglas de firewall, acceso remoto, despliegue de aplicaciones y automatización de configuraciones. El tercero introduce monitoreo, búsqueda y análisis de eventos mediante pipelines de telemetría, logs, Elasticsearch, Kibana y datasets de seguridad. El cuarto integra tecnologías emergentes, particularmente inteligencia artificial aplicada al análisis de datos y blockchain como mecanismo para estudiar integridad, trazabilidad y consenso.

\section{Políticas del curso}
\label{sec:politicas}

A continuación se presentan las políticas generales del curso. Estas políticas son fundamentales para el buen desarrollo del curso y para garantizar un ambiente de aprendizaje efectivo y respetuoso para todos los participantes.

\begin{itemize}
  \item Los estudiantes deben revisar regularmente el correo institucional y el espacio del curso en en la plataforma educativa utilizada. Toda comunicación oficial entre estudiantes y el profesor se realizará a través de estos medios.
  \item Google Classroom es la plataforma única para la gestión del curso en esta edición. En esta plataforma se compartirán avisos urgentes, se controlará la asistencia y se registrarán las evaluaciones. Los estudiantes deben revisar y cargar las actividades, tareas y exámenes asignados en ella. Es responsabilidad de cada estudiante mantenerse informado sobre fechas y actividades. Las actividades entregadas por otros medios o fuera del tiempo estipulado no serán aceptadas ni evaluadas.
  \item La asistencia se registrará en cada sesión. Para tener derecho a presentar los exámenes, los estudiantes deben cumplir con la asistencia mínima establecida en el reglamento escolar.
  \item No se permitirán extensiones en las fechas de entrega de actividades. Cada actividad tendrá un plazo de entrega razonable y bien definido partir de su fecha de asignación que el profesor comunicará en cada asignación. La plataforma educativa estará configurada con las fechas de entrega correspondientes. En caso de cualquier contradicción, prevalecerá la indicación que indique el profesor en la clase.
  \item Las actividades evaluativas indicadas en clase son suficientes para obtener una evaluación adecuada si se realizan en tiempo y forma. No habrá oportunidades adicionales para mejorar evaluaciones después de su entrega. La evaluación final se basará únicamente en las actividades entregadas en los plazos establecidos.
  \item Todas las entregas de actividades deben ser en formato PDF. No se recibirán ni evaluarán actividades entregadas en otros formatos. Si se requiere enviar más de un archivo, estos deben ser entregados en formato PDF y comprimidos en un archivo .ZIP.
  \item Aunque las actividades puedan realizarse en equipo, cada miembro debe enviar su propia copia del trabajo, independientemente de que otro compañero haya enviado la suya. Cada estudiante debe subir el trabajo individualmente en la plataforma del curso. No se aceptará de ninguna forma la justificación “no lo entregué porque mi compañero ya lo entregó”.
  \item Las clases serán en línea serán y en tiempo real, lo que significa que los estudiantes deben participar en el momento programado para la clase. La asistencia en estas sesiones en línea será obligatoria. Es obligatoria la participación con cámaras encendidas y micrófonos apagados, salvo que se indique lo contrario. Los estudiantes que deseen participar activamente deberán hacerlo mediante las opciones disponibles en la plataforma que se utilice.
  \item Si alguna clase debe ser cancelada, el profesor notificará con suficiente antelación. Las clases y actividades canceladas serán reprogramadas para la semana siguiente. La reposición se considerará válida solo si al menos el 85\% de los estudiantes asisten a la clase reprogramada.
  \item Se permitirá una tolerancia de hasta 10 minutos para el ingreso a las clases. Luego de este tiempo no se permitirá la entrada al salón.
  \item El tiempo de descanso durante las clases será acordado con los estudiantes al inicio de cada sesión. El descanso podrá realizarse durante la clase o al final, según lo pactado.
  \item Los estudiantes deben cumplir estrictamente con los reglamentos establecidos, en particular en lo que respecta a comportamiento, respeto y otras cuestiones académicas.
  \item Queda terminantemente prohibido ingresar a las clases con alimentos o bebidas y comer durante las sesiones. Esta política debe ser respetada en todas las clases presenciales.

\end{itemize}

\subsection{Comunicación}
\label{sec:comuinicacion}
Toda la comunicación se hará por los canales oficiales (correo electrónico, plataforma del curso, etc.). La dirección de correo del profesor es: lbustio@inaoe.mx. Las asesorías podrán ser pactadas por los estudiantes y el profesor en un espacio de tiempo de mutua conveniencia.


\subsection{Cronograma}
\label{sec:cronograma}

A continuación se presenta el cronograma tentativo del curso. Este cronograma es una guía general y puede estar sujeto a cambios según las necesidades del curso.

\setlength{\LTleft}{0pt}
\setlength{\LTright}{0pt}
\begin{longtable}{|
  >{\centering\arraybackslash}p{0.04\textwidth}|
  >{\centering\arraybackslash}p{0.07\textwidth}|
  >{\raggedright\arraybackslash}p{0.39\textwidth}|
  >{\raggedright\arraybackslash}p{0.39\textwidth}|}
\caption{Calendario del curso}
\label{tab:calendario-curso}
\\
\hline
\textbf{No} & \textbf{Fecha} & \textbf{Tema} & \textbf{Comentario} \\
\hline
\endfirsthead

\hline
\textbf{No} & \textbf{Fecha} & \textbf{Tema} & \textbf{Comentario} \\
\hline
\endhead

\hline
\endfoot

1 & 11/05 & \textbf{Introduccion al curso}
            \begin{itemize}
              \item Presentacion del curso.
              \item Explicacion de los temas.
              \item Presentacion de la plataforma a emplear en el curso.
            \end{itemize}
&
\textbf{Tarea}
\begin{itemize}
  \item Crear una cuenta de usuario en Google Cloud Platform.
  \item Revisar los videos introductorios [1], [2] y [3].
\end{itemize}
\\
\hline

2 & 15/05 & \textbf{Cloud computing con GCP}
            \begin{itemize}
              \item Enfoque practico, trabajo incremental, documentacion y reglas.
              \item GCP como plataforma del curso.
              \item Cloud computing, modelo cloud, IaaS, PaaS y SaaS.
              \item Modelos de despliegue: nube publica, privada, hibrida y multi-cloud.
              \item Proyectos, Compute Engine y VM Debian.
              \item Regiones, zonas, VPC, firewall e IP publica/privada.
              \item SSH a instancias Linux.
              \item Costos, billing y buenas practicas de ahorro.
              \item Normativas y compliance: GDPR, ISO 27017 e ISO 27018.
              \item Riesgos y troubleshooting de VM, SSH, firewall y web.
              \item Almacenamiento cloud con buckets.
            \end{itemize}
&
\textbf{Tarea}
  \begin{itemize}
    \item Revisar el material de estudio de la sesion.
    \item Crear cuenta/proyecto en GCP.
    \item Desplegar una VM Debian 12 con configuracion minima definida.
    \item Documentar region, zona, costo y buenas practicas de ahorro.
    \item Conectarse por SSH e inspeccionar SO, CPU, memoria y disco.
    \item Instalar herramientas complementarias: \texttt{htop} y \texttt{mc}.
    \item Configurar reglas de firewall HTTP.
    \item Instalar Nginx.
    \item Verificar Nginx desde la VM y desde navegador externo.
    \item Crear bucket en Cloud Storage y subir archivo de prueba.
    \item Entregar reporte con evidencias, diagrama, costos, riesgos, compliance, troubleshooting y reflexion.
  \end{itemize}
\\
\hline
3  & 18/05 & \textbf{Clase práctica: Infraestructura Cloud en GCP}
            \begin{itemize}
              \item Espacio de trabajo para completar la practica asignada en la sesion anterior...
              \item Revision de objetivos tecnicos: proyecto en GCP, VM Debian, acceso SSH, reglas de firewall, Nginx y bucket de Cloud Storage.
              \item Verificacion de evidencias minimas: capturas, comandos ejecutados, configuraciones, pruebas de acceso y resultados observados.
              \item Identificacion y documentacion de problemas tecnicos: errores de SSH, firewall, permisos, servicios, costos y disponibilidad.
              \item Discusion de buenas practicas de seguridad, ahorro de costos y limpieza de recursos.
              \item Preparacion del reporte tecnico individual.
            \end{itemize}
&
  \textbf{Tarea}
  \begin{itemize}
    \item Completar o corregir la practica iniciada en la sesion anterior.
    \item Validar que la VM sea accesible por SSH y que Nginx responda desde navegador externo.
    \item Confirmar la existencia del bucket y la carga de un archivo de prueba.
    \item Registrar comandos relevantes, decisiones tecnicas y problemas encontrados.
    \item Incluir evidencias de funcionamiento, costos estimados, riesgos y reflexion final.
    \item Crear un reporte en PDF que describa la práctica, docuemntando los procesos y resultados.
  \end{itemize}
\\
\hline
4  & 22/05 & \textbf{ELK Stack}
  \begin{itemize}
    \item Introduccion a pipelines modernos de datos: ingestion, indexacion, busqueda y visualizacion.
    \item Elasticsearch: indices, documentos, mappings, consultas y API REST.
    \item Kibana: dashboards, filtros, visualizaciones y exploracion de eventos.
    \item JSON, \texttt{curl} y scripts de ingestion.
    \item Logs, metricas, eventos, telemetria, monitoreo y relacion con ciberseguridad.
    \item Casos de uso: busqueda de eventos, deteccion de actividad sospechosa y analisis temporal.
  \end{itemize}
&
\textbf{Tarea}
  \begin{itemize}
      \item Revisar el material de la clase.
      \item Analizar la arquitectura general de un pipeline ELK.
      \item Identificar componentes, puertos, flujo de datos y requisitos de la VM.
      \item Revisar ejemplos de documentos JSON y consultas REST.
      \item Preparar dudas tecnicas de la siguiente sesión.
      \item Sin tarea entregable en esta sesion.
  \end{itemize}
\\
\hline
5  & 25/05 & \textbf{Instalacion de ELK Stack}
  \begin{itemize}
    \item Preparacion de la VM Linux en GCP.
    \item Requisitos de sistema, memoria, servicios y paquetes.
    \item Instalacion de Elasticsearch.
    \item Instalacion de Kibana.
    \item Configuracion de puertos, firewall y acceso remoto.
    \item Verificacion inicial de servicios.
    \item Troubleshooting de instalacion.
  \end{itemize}
&
\textbf{Tarea}
  \begin{itemize}
    \item Preparar o reutilizar una VM Linux en GCP.
    \item Instalar Elasticsearch.
    \item Instalar Kibana.
    \item Validar estado de los servicios.
    \item Documentar comandos, configuraciones y errores encontrados.
    \item Dejar la VM preparada para la implementacion de la sesion 6.
  \end{itemize}
\\
\hline
6  & 29/05 & \textbf{Detección de Ataques de Fuerza Bruta mediante Análisis Estructurado de Logs en Kibana}
  \begin{itemize}
    \item Puesta en marcha del pipeline de eventos.
    \item Ingestion de eventos o logs en formato JSON.
    \item Indexacion y validacion de documentos en Elasticsearch.
    \item Consultas basicas mediante API REST.
    \item Visualizacion de datos en Kibana.
    \item Construccion de dashboards iniciales.
    \item Analisis de eventos y patrones observados.
  \end{itemize}
&
\textbf{Tarea}
  \begin{itemize}
    \item Indexar un dataset de eventos o logs en formato JSON.
    \item Realizar consultas basicas mediante API REST.
    \item Construir visualizaciones y un dashboard inicial en Kibana.
    \item Analizar eventos relevantes del dataset.
    \item Entregar un reporte donde se describa la arquitectura desarrollada, dataset, consultas, dashboard, problemas encontrados, soluciones y evidencias.
  \end{itemize}
\\
\hline
7  & 01/06 & \textbf{Clase práctica: Detección de Ataques de Fuerza Bruta mediante Análisis Estructurado de Logs en Kibana}
  \begin{itemize}
    \item Sesion de trabajo para consolidar lo realizado en las sesiones del 25/05 y 29/05.
    \item Verificacion de la VM, recursos disponibles, puertos, reglas de firewall y acceso a servicios.
    \item Revision del estado de Elasticsearch y Kibana.
    \item Correccion de errores de instalacion, configuracion, memoria, permisos, conectividad y servicios.
    \item Integracion del flujo completo: ingestion de datos, creacion de indices, consultas REST y visualizacion en Kibana.
    \item Validacion de evidencias tecnicas para documentar el funcionamiento del pipeline.
    \item Cierre integrado de los modulos 1, 2 y 3: infraestructura cloud, pipelines ELK y analisis de seguridad con telemetria.
  \end{itemize}
&
\textbf{Tarea}
  \begin{itemize}
    \item Completar la instalacion funcional de Elasticsearch y Kibana.
    \item Ejecutar la ingestion de eventos o logs en formato JSON.
    \item Validar indices, documentos y consultas basicas en Elasticsearch.
    \item Construir o mejorar visualizaciones y dashboard en Kibana.
    \item Documentar problemas encontrados, soluciones aplicadas, riesgos, costos y decisiones tecnicas.
    \item Entregar un reporte en PDF con la descripcion, definición e implementación de la arquitectura, procedimiento, evidencias, consultas, dashboard, analisis de resultados y reflexion final.
    \item El reporte funciona como evidencia de cierre de los modulos 1, 2 y 3.
  \end{itemize}
\\
\hline
8  & 05/06 & \textbf{Aprendizaje Automatizado en Google Cloud Platform}
  \begin{itemize}
    \item Introduccion practica al flujo completo de Machine Learning en GCP.
    \item Uso de Cloud Storage para almacenar datasets y artefactos de modelo.
    \item Creacion de un notebook en Colab Enterprise como entorno de desarrollo administrado.
    \item Carga, exploracion y preparacion del dataset Titanic.
    \item Entrenamiento de un modelo supervisado con Python, Pandas y Scikit-Learn.
    \item Evaluacion del modelo mediante accuracy, matriz de confusion e importancia de variables.
    \item Serializacion del modelo entrenado y almacenamiento en Cloud Storage.
    \item Despliegue del modelo como API REST mediante Cloud Run.
    \item Prueba del servicio de inferencia desde un cliente externo mediante \texttt{curl}.
  \end{itemize}
&
\textbf{Tarea}
  \begin{itemize}
    \item Crear o reutilizar un bucket de Cloud Storage con la estructura \texttt{data/} y \texttt{models/}.
    \item Subir el dataset \texttt{titanic.csv} y cargarlo desde Colab Enterprise.
    \item Ejecutar el proceso completo de preparacion de datos, entrenamiento, evaluacion y prediccion.
    \item Guardar el modelo entrenado como \texttt{titanic\_model.pkl} y subirlo a Cloud Storage.
    \item Construir los archivos \texttt{app.py}, \texttt{requirements.txt} y \texttt{Dockerfile}.
    \item Desplegar la API en Cloud Run y validar los endpoints de estado y prediccion.
    \item Documentar arquitectura, codigo ejecutado, evidencias, errores encontrados, soluciones aplicadas, resultados del modelo y reflexion sobre el uso de ML en cloud.
    \item Entregar reporte en PDF con capturas, comandos relevantes, URL del servicio o evidencia de despliegue, prueba de inferencia y analisis final.
  \end{itemize}
\\
\hline
9  & 08/06 & & \\
\hline
10 & 12/06 & & \\
\hline
11 & 15/06 & & \\
\hline
12 & 19/06 & & \\
\hline
13 & 22/06 & & \\
\hline
14 & 26/06 & & \\
\hline
15 & 29/06 & & \\
\hline
16 & 03/07 & & \\
\hline

\end{longtable}

\noindent\textbf{Recursos}

\begin{itemize}
  \item[{[1]}] \url{https://www.youtube.com/watch?v=lvZk_sc8u5I}
  \item[{[2]}] \url{https://www.youtube.com/watch?v=4dNSAIwXO5M}
  \item[{[3]}] \url{https://www.youtube.com/watch?v=HU58N5fz7B8}
\end{itemize}


\subsection{Integridad}
Se prohíbe presentar como propio material ajeno, reportes generados sin verificación o evidencias que no correspondan a prácticas realizadas por el estudiante. La integridad académica implica documentar fuentes, describir procesos reales, reconocer errores técnicos y reportar resultados verificables.

\section{Metodología}
La metodología del curso combina explicación conceptual, demostración técnica, práctica guiada, trabajo autónomo y discusión crítica. Cada actividad busca que el estudiante implemente infraestructura o servicios reales, documente decisiones técnicas, identifique riesgos y analice problemas encontrados durante la configuración.

Se fomentará el uso de una bitácora técnica. La bitácora debe incluir fecha, objetivo de la práctica, servicios utilizados, comandos relevantes, configuraciones aplicadas, evidencias, problemas encontrados, soluciones intentadas, costos observados y reflexión final.

\section{Evaluación}
La evaluación se compone de prácticas, reportes técnicos, discusión de resultados y proyecto final.

\begin{center}
\begin{tabular}{l r}
\hline
Componente & Porcentaje \\
\hline
Prácticas de laboratorio & 40\% \\
Reportes técnicos & 30\% \\
Proyecto final & 30\% \\
\hline
\end{tabular}
\end{center}

\subsection{Criterios de prácticas}
Las prácticas se evaluarán con una rúbrica que incluye:
\begin{itemize}
  \item Infraestructura funcional y verificable.
  \item Configuración técnica coherente con los objetivos de la actividad.
  \item Evidencias claras de funcionamiento.
  \item Documentación reproducible del procedimiento.
  \item Identificación de riesgos, costos y problemas técnicos.
  \item Análisis crítico de los resultados obtenidos.
\end{itemize}

\section{Proyecto final}
El proyecto final consistirá en diseñar, implementar y documentar una solución apoyada en infraestructura cloud y tecnologías emergentes revisadas durante el curso. El estudiante deberá justificar la arquitectura, describir los servicios utilizados, presentar evidencias de funcionamiento, analizar riesgos y discutir limitaciones técnicas.

\section{Bibliografía}

La bibliografía integra los textos base del curso, los libros disponibles en la carpeta \texttt{07.- bibliografía}, el repositorio utilizado para el despliegue de ELK Stack y los materiales de apoyo empleados en los temas individuales. No se incluyen materiales ubicados en carpetas denominadas Complementario o Complementarios.

\begin{sloppypar}
\begin{itemize}[leftmargin=*]
  \item Antonopoulos, A. M., \& Wood, G. (2019). \textit{Mastering Ethereum: Building smart contracts and DApps}. O'Reilly Media.

  \item Bustio-Martínez, L. (2026). \textit{Introducción al curso} [Material de clase]. Instituto Nacional de Astrofísica, Óptica y Electrónica.

  \item Bustio-Martínez, L. (2026). \textit{Google Cloud Platform} [Material de clase]. Instituto Nacional de Astrofísica, Óptica y Electrónica.

  \item Bustio-Martínez, L. (2026). \textit{ELK Stack} [Material de clase]. Instituto Nacional de Astrofísica, Óptica y Electrónica.

  \item Bustio-Martínez, L. (2026). \textit{Instalación de ELK Stack} [Laboratorio]. Instituto Nacional de Astrofísica, Óptica y Electrónica.

  \item deviantony. (s. f.). \textit{Elastic stack (ELK) on Docker}. GitHub. Recuperado el 23 de mayo de 2026, de \url{https://github.com/deviantony/docker-elk}

  \item Dixit, B., Kuć, R., Rogoziński, M., \& Chhajed, S. (2017). \textit{Elasticsearch: A complete guide}. Packt Publishing.

  \item Foster, I., \& Gannon, D. B. (2017). \textit{Cloud computing for science and engineering}. MIT Press.

  \item Google Cloud. (s. f.). \textit{Google Cloud Fundamentals: Core Infrastructure (en español)}. Coursera. \url{https://www.coursera.org/learn/gcp-fundamentals-es}

  \item Hwang, K. (2017). \textit{Cloud computing for machine learning and cognitive applications}. MIT Press.

  \item Vacca, J. R. (Ed.). (2016). \textit{Cloud computing security: Foundations and challenges}. CRC Press.

  \item Widjaja, A. (2022). \textit{Data engineering with Google Cloud Platform: A practical guide to operationalizing scalable data analytics systems on GCP}. Packt Publishing.

  \item YouTube. (s. f.). \textit{Introducción a Google Cloud Platform}. \url{https://www.youtube.com/watch?v=lvZk_sc8u5I}

  \item YouTube. (s. f.). \textit{Google Cloud Platform Tutorial}. \url{https://www.youtube.com/watch?v=4dNSAIwXO5M}

  \item YouTube. (s. f.). \textit{Google Cloud Platform for Beginners}. \url{https://www.youtube.com/watch?v=HU58N5fz7B8}
\end{itemize}
\end{sloppypar}

\section{Cierre}
Esta guía sintetiza la estructura general del curso, sus políticas, el cronograma tentativo y los criterios de evaluación. Las actividades específicas podrán ajustarse según el avance del grupo, disponibilidad de plataformas y necesidades técnicas del curso.

\end{document}


